% Table S12: Target-specific RRS summaries
\begin{table}[htbp]
\centering
\scriptsize
\caption{Target-specific docking-score retention in the eligible candidates. Values are summarized across candidates with an eligible WT denominator for the indicated target. The target-specific class counts are descriptive and are not equivalent to the primary two-target classification.}
\label{tab:s12_rrs_by_target}
\begin{tabular}{l c c c c c c}
\toprule
Target & Eligible candidates & Mean RRS & Median RRS & Range & Class A*/A & Class B/C/D \\
\midrule
PfDHFR & 12 & 75.1 & 70.8 & 59.7--123.4 & 2 & 10 \\
PfCRT & 17 & 86.2 & 83.5 & 74.7--104.8 & 10 & 7 \\
\bottomrule
\end{tabular}
\end{table}

\noindent\textit{State-specific summaries.} For PfDHFR, mean RRS values were 74.7 (N51I), 73.7 (C59R), 75.3 (S108N), and 76.8 (I164L). For PfCRT, the corresponding means were 85.4 (K76T) and 87.0 (K76A). These target-specific summaries describe score retention and do not measure biochemical resistance.
